\chapter{脂质：不只是“油和肥肉”}

\begin{kaichang}
做一个小观察。倒半杯清水，加一勺食用油，用筷子拼命搅三十秒。搅拌的时候，油被扯成无数金黄的小珠，在水里四散翻滚，看起来像是要“化”进去了。可只要你停手，不出一分钟，小珠子们就一颗颗找回同伴，重新聚成完整的一层，浮在水面上，界限分明，仿佛刚才的搅拌从没发生过。

再看厨房里的另两个现象。冬天，那碗猪油是乳白色的固体，用勺子挖都费劲；旁边那瓶菜籽油却清亮流动。同样是“脂肪”，凭什么一个凝固、一个流动？还有蛋黄——做菜的人都知道，蛋黄能把油和水“撮合”在一起做成蛋黄酱，搅多久都不分层。它凭什么做到水和油自己都做不到的事？

最后一个问题留给你身上的肥肉：它和锅里的食用油，是不是同一种东西？它除了“让你长胖”，还在替你干哪些你从没感谢过它的活？

这一章结束时，这几个问题会串成一条完整的因果链。先把你对油水分层的猜测写下来。
\end{kaichang}

\section{看得见的脂肪：不溶、浮起、会凝固}

先把“能观察到什么”说清楚，再谈解释。

油和水不相混，这你已经亲眼看到了。它浮在水面而不是沉底，说明油的密度比水略小（大约在水的九成上下）。油摸上去滑腻，不容易挥发——一杯水放几天会干掉，一杯油放几天几乎不少。油能烧：锅里油温太高会起火，所以油锅起火\textbf{绝不能用水浇}（水沉到油下剧烈汽化，把燃烧的油滴炸得到处都是），正确做法是盖上锅盖隔绝空气。

再看凝固。猪油、黄油在室温附近是软的固体，花生油、菜籽油是液体，而它们加热后都会变成液体，冷却后各回各的状态。这说明“固还是液”不是油和脂肪的本质区别，只是熔点落在了室温的哪一边。顺便记住一个术语：化学上把常温下呈液态的叫“油”、呈固态的叫“脂”，但它们属于同一个大家族——\textbf{脂质}。

还有一类现象容易被忽略：牛奶是均匀的乳白色，放多久也不会“油水分家”（只会微微浮起一层奶油）；蛋黄酱更是油和水的稳定混合。可见“油水必定分层”这条经验，在某些物质在场时会失效。厨房里早就藏着例外，只是你没拿它当回事。

\section{把镜头推进：一条尾巴和一颗头}

为什么脂肪和水合不来？把镜头推进到分子层。这一家子的成员不少，先认识四位主角。

\textbf{第一位：脂肪酸。}它是理解一切的起点。一个脂肪酸分子可以分成两段：一条十几个碳原子连成的长链，链上挂满氢——这一段纯粹是碳和氢，和汽油、石蜡、聚乙烯是同一种“脾气”（第 36 章说过，碳氢链排斥水、亲近油）；链的尽头是一个羧基（\chem{-COOH}，第 37 章的官能团），带氧、能和水分子握手。一个分子，拖着一条“亲油的长尾巴”，顶着一颗“亲水的小头”。这个“一头两脾气”的结构，是本章所有故事的种子。

\textbf{第二位：甘油三酯。}单个脂肪酸不好存放，细胞把三个脂肪酸分子的羧基“挂”到一个叫甘油的小分子骨架上，通过酯键连接（酯化反应，第 40 章已经见过），得到的大家伙就是甘油三酯。你吃的食用油、猪油、你身上的肥肉，主要成分都是它。三个羧基全被占用了，于是甘油三酯\textbf{几乎不带亲水的头}——它是一条纯粹的三尾章鱼，彻头彻尾地“不合群”。这才是“油”的真身。

\textbf{第三位：磷脂。}把甘油三酯的其中一条脂肪酸尾巴换掉，接上一个带磷酸基团的大号亲水的头，就得到磷脂。它比脂肪酸更极端：两条粗壮的亲油尾巴，配一个又大又强势的头。蛋黄里的卵磷脂、你每一个细胞的膜，都是它。磷脂是“撮合油水”的专业选手——蛋黄酱的秘密就在它身上。

\textbf{第四位：胆固醇。}它和前三位长得完全不像：没有长尾巴，主体是四个碳环拼成的一块硬邦邦的“小板”，板上只带一个小小的羟基（\chem{-OH}）。它是细胞膜的“填缝料”和许多激素的原料，名声很坏，用处很大——后面专门为它辟一节。

四位主角，三种脾气：纯亲油（甘油三酯）、一头一尾（脂肪酸、磷脂）、小板一块（胆固醇）。记住这个分类，后面的规则全从脾气里长出来。

\section{尾巴的形状：饱和、不饱和与反式}

现在回答开场问题之二：凭什么猪油是固体、菜籽油是液体？

差别不在“动物油还是植物油”这个出身，在脂肪酸尾巴的\textbf{形状}。第 36 章讲过碳骨架的两条关节规则：单键能自由转动，双键不能转动。如果一条碳链上全是单键，链可以舒展得笔直，化学上叫\textbf{饱和}脂肪酸（碳上能挂的氢挂满了，“饱和”）；如果链中间有一个碳碳双键，双键像焊死的关节，把链固定出一个弯——这样的叫\textbf{不饱和}脂肪酸。

形状决定堆积，堆积决定熔点。笔直的尾巴能像筷子一样整齐码放，分子之间靠得近、接触面积大，第 22 章讲的分子间作用力攒起来就很可观，需要更高的温度才能把它们抖散——所以动物脂肪里占多数的饱和脂肪酸熔点高，室温下是固体。弯了的尾巴码不整齐，像一筐弯吸管，松松垮垮，稍一升温就散开——植物油里占多数的不饱和脂肪酸熔点低，室温下是液体。“油还是脂”，说到底是\textbf{尾巴弯不弯}的问题。

还有一个细节值得说：双键锁死后，链的两段可以弯在同侧（顺式），也可以分处两侧（反式）。天然植物油里的双键几乎全是顺式，弯得厉害。而工业上把植物油部分加氢（往双键上补氢，让油变稠变耐放，做人造奶油、起酥油）时，会有一部分双键没被补氢、却被“掰直”成了反式——分子变直了，熔点升高了，口感变好了，货架期变长了，但它在你体内的表现比顺式更糟：大量人群研究显示，反式脂肪摄入与心血管疾病风险升高相关，而且看不到什么好处，于是许多国家已经限制或禁用工业生产的反式脂肪。这里要注意证据的措辞：是“风险升高相关”，不是“碰一口就中毒”——剂量永远重要，偶尔吃到和长期大量吃是两回事。看配料表时，“氢化植物油”“人造奶油”“起酥油”是旧工艺的线索，不过配方在改进，具体含量要看营养成分表。

\begin{qiaoliang}
为什么双键能把熔点压低几十度？算一笔接触的账。把脂肪酸尾巴简化成一排珠子：两条笔直的 18 碳链并排贴紧，大约有十几对“珠子挨着珠子”的接触点，每个接触点贡献一份微弱的分子间作用力。中间弯一个顺式双键，两条链像错开的拉链，接触点可能只剩一半不到。分子间作用力每份都很小，但熔点比的就是这些微小吸引力的\textbf{总和}对抗热运动。总和砍掉一半，需要的“抖动温度”自然大降——硬脂酸（全直，18 碳）熔点约 $69\,^{\circ}\mathrm{C}$，同样 18 碳、只多一个顺式双键的油酸，熔点骤降到约 $13\,^{\circ}\mathrm{C}$。一个双键，五十多度的差别——这就是“结构决定性质”最直白的一次演示。
\end{qiaoliang}

\section{不是油“怕”水，是水太爱自己}

接下来是全章最重要的一条规则。我们通常说“油疏水”，听起来像油分子讨厌水分子、躲着走。错了。分子没有好恶，而且油和水之间其实也存在微弱的吸引力。真正的主角是\textbf{水和水自己}。

回忆第 21、22 章：水分子之间靠氢键互相拉扯，形成一张不断断裂又重组的大网。这张网对水来说能量上很划算——每个水分子都想方设法多握住几只邻居的“手”。现在，一个油分子掉进水里：它不带电、不能形成氢键，像一颗挤进舞会的石头。它周围的水分子够不着它，只好彼此抱团，在油分子表面排成一个比平常更有序的“笼子”，白白损失了一批本来可以自由组合的氢键搭档。从统计上看（第 28 章的语言），这种“笼子水”的排法数目少、熵低，不合算。

怎么办？系统有一个极其简单的出路：\textbf{把油分子们赶到一起去}。一千个油分子如果四散分布，水要给它们造一千个小笼子；如果它们聚成一个大油球，笼子只需要包住球面，被束缚的水分子少得多，重获自由的水分子多得多——总熵上升。于是你搅拌时扯散的油珠，不是“想家”了才团聚，而是被水分子集体的统计账“挤”回了同伴身边。这就是\textbf{疏水效应}：它不是油和水之间的一种新力，而是水的氢键网络为了少受损失，把不速之客“外包”出去的结果。

现在轮到磷脂表演了。磷脂只有一颗亲水的头，尾巴却有两条，所以它既不能整个溶于水，也不甘心整个被挤成油球。它的折中方案漂亮极了：把头伸进水里，把尾巴藏进内部，成千上万分子自动排成\textbf{两层}——头朝外面对水，尾巴在夹层里互相抱团。不需要图纸，不需要监工，把磷脂撒进水里，它自己就会组装成双层，还会弯成封闭的小球，把一小包水裹在里面。

\begin{center}
\begin{tikzpicture}[scale=1]
\foreach \x in {0,1,2,3,4,5} {
  \fill[blue!50] (\x,2.0) circle (0.15);
  \draw[thick] (\x,1.85) -- (\x-0.08,1.1);
  \draw[thick] (\x,1.85) -- (\x+0.08,1.1);
  \fill[blue!50] (\x,0.2) circle (0.15);
  \draw[thick] (\x,0.35) -- (\x-0.08,1.1);
  \draw[thick] (\x,0.35) -- (\x+0.08,1.1);
}
\node at (2.5,2.6) {水};
\node at (2.5,-0.4) {水};
\node at (7.3,1.1) {尾巴躲进夹层};
\end{tikzpicture}
\end{center}

图里圆圈是亲水的头，短线是亲油的尾巴——这是人为的表示法，真实的磷脂分子没有这么圆的头和这么直的尾巴。两层磷脂组成的“膜”，厚度只有约 5 纳米，却是每一个细胞的外墙。它的建造方式震撼就震撼在：\textbf{没有建筑商，只有疏水效应这一条统计规律}。膜的完整故事——它怎么控制物质进出、怎么产生电压——是第 51 章的内容，这里你只要先相信：边界可以自己长出来。

\begin{shiyan}
\textbf{让油和水握一次手：乳化实验}

材料：三个一样的透明小瓶（带盖）、食用油、水、一个生鸡蛋黄、一滴洗洁精、一根牙签。

步骤：
\begin{enumerate}[itemsep=1pt]
\item 三个瓶子各装约三分之一水，再各加一勺食用油。
\item 1 号瓶什么都不加，盖紧摇 20 秒，静置，计时观察分层。
\item 2 号瓶用牙签挑一点蛋黄进去，同样摇 20 秒，静置观察。
\item 3 号瓶加一滴洗洁精，摇 20 秒，静置观察。
\end{enumerate}

观察点：1 号瓶是不是一分钟内就界限分明？2 号瓶的混合物是不是变得浑黄浓稠、很久不分层？3 号瓶的泡沫和浑浊能维持多久？想一想：蛋黄里的卵磷脂和洗洁精分子，站在油和水的交界面上时，头和尾分别朝着哪边？

安全提示：只用家常物品，可以放心做；但实验用过的材料\textbf{不要再吃}；生鸡蛋可能带细菌，做完洗手；瓶子洗净再回收。
\end{shiyan}

\section{符号登场：把四位主角写出来}

概念清楚了，符号才有意义。用化学家的写法把四位主角登记在册。

硬脂酸，一条 18 碳的直链饱和脂肪酸：\[\chem{CH_3(CH_2)_{16}COOH}\]一共 17 个碳在链上，加羧基里那 1 个，正好 18 个。油酸，同样 18 碳，只在正中间多一个顺式双键：\[\chem{CH_3(CH_2)_7CH{=}CH(CH_2)_7COOH}\]看出规律了吗？碳数相同，只差一个双键——就是这一个焊死的弯，把熔点压低了五十多度。

甘油三酯的形成，是第 38 章讲过的缩合反应的一次三连发：一个甘油分子有三个羟基（\chem{-OH}），分别和三个脂肪酸的羧基各挤出一分子水，连成三个酯键：\[\chem{C_3H_8O_3} + 3\,\chem{C_{17}H_{35}COOH} \rightarrow \chem{C_{57}H_{110}O_6} + 3\,\chem{H_2O}\]产物就是一个典型的动物脂肪分子（这里三个脂肪酸都写成硬脂酸；真实油脂里三条尾巴通常不一样长、弯直混搭）。反过来，加水把酯键剪开就是水解——你的消化系统靠它拆油，肥皂厂靠它做肥皂（第 40 章的皂化反应），同一根键，两种拆法。

胆固醇的完整结构式比较复杂，你只需要记住它的骨架是四个稠合的碳环（三个六元环加一个五元环），分子式 \chem{C_{27}H_{46}O}——27 个碳里几乎全是碳氢，只有一个氧，所以它和油一样怕水，绝不单独在血液里游荡。这个性质马上有大用场。

\section{一克脂肪的能量账}

肥肉除了让你坐下去时多一层软垫，到底在干什么？先算它最著名的那笔账：储能。

\begin{qiaoliang}
同样烧掉 1 克，脂肪释放约 \ud{38}{kJ}，糖只有约 \ud{17}{kJ}——差了两倍多。为什么？关键在“已经烧过没有”。燃烧的本质是碳和氢与氧化合（第 26 章）。糖的骨架上挂满了氧（葡萄糖 \chem{C_6H_{12}O_6}，碳和氧几乎一样多），相当于碳已经“烧掉一半”了，再烧放出的能量自然有限；而脂肪酸的碳链上几乎全是碳氢键，每一个碳都还是“未烧的新柴”，从头到尾氧化一遍，放能自然多。

更要命的是储存成本。细胞存糖用的是糖原，而糖原亲水，每存 1 克糖原要绑走 2～3 克水；脂肪疏水，存起来就是干干净净的纯能量。算上水的重量，同样重量下脂肪的实际储能密度能达到糖的六倍以上。对一头要迁徙、要冬眠、要挨饿的动物来说，用糖储能等于背着水箱赶路，用脂肪才是真轻装。
\end{qiaoliang}

算一个你身上的数量级。一个普通成年人身上约有 \ud{10}{kg} 脂肪（胖瘦不同，差别很大），按每克 \ud{38}{kJ} 算是 $10\,000 \times 38 \approx \ud{3.8\times 10^{5}}{kJ}$。一个人一天总共消耗约 \ud{8000}{kJ}。两数相除：这份储备理论上够烧四十多天。而全身的糖原储备只有几百克，连一天都不够。候鸟能一口气飞越几千公里不进食，骆驼的驼峰里塞的也不是水而是十几公斤脂肪（脂肪氧化时还会产生水，顺便补水），靠的都是这笔账。当然，这是理想化的估算——真实饥饿时身体会同时分解蛋白质，激素也会层层调节，人是烧不“干净”这四十天的；但数量级本身已经说明了进化的选择：长距离、长周期的能量储备，非脂肪不可。

\section{保温、缓冲，还有信号}

储能只是肥肉的简历第一页。

\textbf{保温。}脂肪导热慢，是贴身的隔热层。海豹、鲸、企鹅在冰水里活动自如，靠的就是皮下那层几厘米厚的脂肪；你冬天比别人更怕冷或更抗冻，一部分原因也写在皮下脂肪厚度里。鲸类的这层脂肪有个专门的名字叫鲸脂，既保温又储能，一举两得。

\textbf{缓冲。}你的肾脏、眼球都不是直接顶在体壁上，而是被柔软的脂肪垫托着、裹着——走路、跳跃、磕碰时，这层垫子吸收冲击。没有它，一次普通的摔跤就可能伤到内脏。

\textbf{信号与原料。}这一栏最容易被忽视，也最颠覆“肥肉是死库存”的印象。胆固醇那块“小板”是一整条生产线的起点：性激素（雌激素、睾酮）、肾上腺皮质激素、帮助消化脂肪的胆汁酸、皮肤晒太阳后合成的维生素 D，全是从胆固醇改造来的。没有胆固醇，这条线全线停产。脂肪细胞本身也不是闷声存货的仓库——它会分泌信号分子（比如瘦素）向大脑汇报库存，参与调节食欲。换句话说，肥肉还是一个内分泌器官，一直在跟你的大脑打电话。

\section{科学家怎样知道：一层膜，还是两层？}

\begin{zhengju}
细胞膜是一层磷脂，还是两层？这个问题听起来没法回答——膜太薄了，上世纪最好的光学显微镜也看不见。1925 年，两位荷兰科学家戈特（Evert Gorter）和格伦德尔（François Grendel）用一个漂亮的“绕道”实验给出了答案。

他们选中了红细胞：成熟的红细胞没有细胞核和细胞器，膜几乎只有最外面那一层，不用操心别的膜来捣乱。第一步，从已知数量的红细胞里把脂质全部提取出来；第二步，把这些脂质滴在水面上，让它们摊开成只有一个分子厚的薄膜；第三步，测量这层薄膜的总面积，和红细胞的表面积比一比。

结果：摊开的面积，大约是红细胞表面积的\textbf{两倍}。如果细胞膜只有一层磷脂，面积应该和细胞表面一样大；是两倍，说明膜是双层——正好对应你今天在教科书里看到的磷脂双层结构。

但这个故事还有后半段，它同样重要。后来的研究者重算发现，戈特他们的实验有两个互相抵消的误差：一来红细胞表面积算小了（细胞其实是扁圆饼，不是他们假设的光滑球），二来用丙酮提取脂质并不完全。两个误差凑巧抵消，答案居然还是对的。二十世纪五十年代电子显微镜问世，人们直接“看见”了膜的结构，双层的结论才真正坐实。这个案例值得记住的教训是：一个经典实验的结论可以是对的，而论证过程却可能是漏的——科学从来不是某个实验一锤定音，而是不同方法互相校对、误差被逐个揪出来的漫长过程。
\end{zhengju}

\section{胆固醇的运输队：脂蛋白}

现在处理全章最容易被说歪的话题。

血液的主要成分是水。而甘油三酯和胆固醇都怕水——它们要在血管这条“运河”里运输，必须坐“船”。船就是\textbf{脂蛋白}：一个球形的颗粒，内部装着甘油三酯和胆固醇，外壳是单层磷脂（头朝外对着血液）加蛋白质，蛋白质还兼任“收货地址标签”。身体造了好几种规格的船，按密度从大到小，最常被谈论的是两种：低密度脂蛋白（LDL）和高密度脂蛋白（HDL）。

于是有了那句尽人皆知的口诀：LDL 是“坏胆固醇”，HDL 是“好胆固醇”。这句话对了一半，错了另外要紧的一半，必须说清楚它的局限：

\begin{itemize}[itemsep=2pt]
\item \textbf{LDL 和 HDL 都不是胆固醇。}它们是船，胆固醇是货，两艘船运的是同一种货。胆固醇分子本身不分好坏——它是细胞膜和激素的必需品，人体每天自己合成约 1 克，比你从食物里吃到的还多。
\item \textbf{“坏”的实质是去向。}大量证据（人群队列、遗传研究和随机对照药物试验层层印证）表明：血液里 LDL 颗粒长期过多，颗粒会钻进动脉血管壁内皮下堆积，启动粥样硬化斑块的形成——方向是明确的因果，不只是相关。所以说 LDL“高”是风险因素，是有硬证据的。
\item \textbf{但“好”的那半边远没有口诀里那么硬。}HDL 负责把组织里多余的胆固醇运回肝脏回收，人群统计里 HDL 偏高的人心血管风险偏低，于是得了“好”的名声。然而“相关”不等于“因果”：后来几种专门升高 HDL 的药物，在大型随机对照试验里\textbf{并没能减少心梗和中风}。这强烈提示 HDL 的高低更可能是身体状态的“仪表盘”，而不是能直接拧动的“方向盘”。
\item \textbf{风险是一个整体，不是一个数。}LDL 还分大小亚型，颗粒数量比总质量更贴近风险；吸烟、血压、血糖、家族史都在同一张风险表上。体检报告上那两个箭头的意义，是医生结合你的整体情况解读的——本章只解释机制，不给任何个体的诊疗建议。
\end{itemize}

一句话总结局限：“好胆固醇、坏胆固醇”把“船”错当成了“货”，把“相关”混同了“因果”，还把一个多维的群体风险压扁成了两个数字的口诀。作为科普缩写它够格，作为行动指南它不够。

\begin{wuqu}
“胆固醇是有害物质，吃得越少越好，最好降到零。”——错在两头。第一，胆固醇是必需品：没有它，细胞膜维持不了正常的流动性，性激素、维生素 D、胆汁酸全部断供；即便你一口胆固醇都不吃，肝脏也会自己开工合成，每天约 1 克。第二，血液指标不是越低越好：风险来自 LDL 颗粒在血管壁的\textbf{长期过量堆积}，而不是胆固醇这种物质本身的存在。真正需要管理的，是运输系统的长期失衡，而不是把一种生命必需的分子当敌人。顺带拆一个相关的广告话术：“本品不含胆固醇”印在植物油瓶上纯属废话——植物本来就不合成胆固醇，这就像在矿泉水瓶上写“不含牛肉”一样正确而多余。
\end{wuqu}

\section{边界：这些模型在哪里会失手}

每章模型都要交代边界，本章有三条。

\textbf{第一，疏水效应不是绝对的力。}它是水的氢键网络给出的统计结果，强弱随温度变化（多数情况下升温会增强它）；遇到既亲水又亲油的两亲分子，它给出的不是“分层”而是各种有序组装——胶束、双层、囊泡，具体组装成什么样，取决于分子头尾的大小比例。把“疏水”当成油水之间有一条鸿沟，就会漏掉这整片自组装的风景。

\textbf{第二，“饱和就固、不饱和就液”是简化。}天然油脂是几十种甘油三酯的混合物，没有单一熔点，只有一段软化的温度范围；而且同样的脂肪还能排出好几种晶体队形，熔点和口感各不相同——第 24 章那盒起白霜的巧克力，就是可可脂换队形的现场。尾巴的弯直只是熔点的第一近似，不是全部。

\textbf{第三，也是最容易越界的一条：膳食脂肪与健康的关系，证据等级参差不齐。}“反式脂肪有害”“LDL 颗粒过多促动脉粥样硬化”有人群、遗传和随机试验的多重印证，属于较硬的结论；而“某种油抗癌”“某种脂肪酸补脑”“吃脂肪直接变肥肉”这类说法，多半只有观察性研究甚至只有营销文案撑腰。观察性研究擅长发现线索，却天生分不清因果——爱吃某种油的人可能同时运动更多、收入更高、看病更勤。剂量、个体差异、整体膳食结构，每一项都能翻转结论。本章能给你的是机制和判断证据的方法，具体到你个人该吃什么，请交给医生和营养师。

\section{回到那杯油水}

开场的问题，现在可以逐条对答案了。

油为什么搅拌之后还要分层？因为水分子之间那张氢键大网不肯为油分子多造笼子——把油挤成一团、界面缩到最小，是整张网的统计账，不是油在“赌气”。油浮在上层，只是密度比水小了那么一点。

猪油为什么是固体、菜籽油为什么是液体？数尾巴上的双键。猪油的脂肪酸尾巴大多笔直，码得整齐，分子间作用力攒得厚，室温抖不散；菜籽油的尾巴带着顺式的弯，堆不严实，室温下早就散成液体。同一家族，一个弯的差别。

蛋黄凭什么撮合油和水？因为卵磷脂一颗头两条尾，站在油和水的交界面上，头扎进水里、尾巴插进油里，把无数小油珠挨个“包裹登记”，水再也挤不动它们——这叫乳化。牛奶出厂前要“均质”，就是把脂肪球打得极小，让蛋白质和磷脂把它们裹住，于是你喝到的牛奶不再漂着油花。

至于你身上的肥肉：它和食用油确实是同一类分子——甘油三酯。但它不是死库存：它是够烧四十天的能量储备，是肾脏和眼球的减震垫，是冬天贴身的隔热层，还是一个不断向大脑发信号的内分泌器官。它出过问题（库存长期过剩会引发炎症和代谢紊乱，那是医学的内容），但它从来不是身体的错误，而是几十亿年进化打磨出的储能方案——在食物没有保证的年代，这套方案是你所有祖先能活到把你生下来的原因之一。

\section{往前一步：边界有了，谁来守门}

这一章我们造出了一种了不起的分子：磷脂。它不需要任何指令，只凭疏水效应就能在水中自动围出一间“单间”——一个把“内”和“外”分开的边界。第 45 章说过，边界是生命的第一件行李；现在你亲眼看到这件行李是怎么自己打包出来的。第 51 章会回到这层膜，细看它怎样选择性地放行和拦截、怎样在两侧维持电压——那是一整套精密的海关系统。

但你可能已经发现了问题：一间全自动围出来的单间，墙壁密不透水，那里面的东西怎么吃饭、怎么扔垃圾？纯磷脂双层对水和离子几乎是封锁的。真正的细胞膜上，嵌满了负责开门、抽水、收发信号的装置——而这些装置，全都不是脂质造的。它们属于下一章的主角：一类用二十种零件折叠出来的分子机器，蛋白质。脂质搭好了舞台，第 48 章，演员登场。

\begin{tiaozhan}
（跳过不影响主线）疏水效应的账其实可以写成一个式子：$\Delta G = \Delta H - T\Delta S$（第 28 章）。油分子分散在水里时，周围“笼子水”排列有序，熵低；油聚成一团后，大批水分子重获自由，$\Delta S$ 为正，于是 $-T\Delta S$ 为负，把 $\Delta G$ 拉向负值——过程自发。注意一个反直觉的细节：这个“油水分离”的驱动力主要来自\textbf{熵}而不是能量——油水混合物的能量（$\Delta H$）变化并不大，有时甚至是吸热的，但熵增照样能把事情办成。更有意思的是，同样的物理规则还有第二份全职工作：蛋白质折叠时，也是疏水的氨基酸侧链被“挤”进分子内部、亲水的留在外面——第 48 章你会看到，蛋白质的三维形状很大程度上是同一条统计规律在水中写下的。
\end{tiaozhan}

\section*{练一练}

\begin{sikaoti}
\item 画一幅磷脂双层的示意图：用圆圈表示亲水的头、折线表示亲油的尾，标出哪一侧接触水，并在图旁注明“这是人为的表示法”。再画一个磷脂在水中围成的封闭小泡，指出小泡内部装的是什么。
\item 算一笔储能账：一名远足者要携带释放约 $\ud{40000}{kJ}$ 能量的储备，分别用脂肪和糖原（糖原每克供能约 $\ud{17}{kJ}$，且每克糖原绑定约 2 克水）需要各背多少公斤？用这个结果向同伴解释：为什么冬眠动物囤脂肪而不是囤糖。
\item 人造奶油厂想把液态植物油变成室温下可涂抹的固体。根据本章规则，他们有两条思路：给双键加氢，或把双键“掰直”。两条路各自改变了分子的什么？各自付出了什么代价（提示：一条失去营养上的优点，一条产生反式脂肪）？
\item 一位长辈说：“体检报告上‘好胆固醇’高，我就可以放心了。”用“船和货”的比喻写三句话，向他解释这句话对在哪里、靠不住在哪里。
\item 戈特和格伦德尔测出“脂质摊开面积约为红细胞表面积的两倍”，于是断定膜是双层。假如当年他们测出的是 1.5 倍，他们能直接宣布“膜是一层半”吗？列出至少两种可能的实验误差来源，并说明你会补做什么测量来分辨它们。
\item 本章说油水分层的驱动力“主要来自熵而不是能量”。用第 28 章的直觉解释：为什么说被赶到一起的油分子并没有“变得更整齐就违反熵增”？（提示：数熵的时候，你只数了油，还是连水一起数了？）
\end{sikaoti}
